% SPDX-License-Identifier: Apache-2.0
%
% The cover: what documents TxHDL has, and which to read for what.
% Built by //docs:cover.
\documentclass[journal,onecolumn,9pt]{IEEEtran}

% One column: 80 columns fit at the body size, so nothing is reduced.
\newcommand{\listingsize}{\normalsize}
% SPDX-License-Identifier: Apache-2.0
%
% The house preamble, shared by every document in this directory. Loaded
% after \documentclass, before \title. Keeping it in one file is what
% stops two articles drifting apart in font, listing style or figure
% style.
% Computer Modern. IEEEtran selects Times (ptm) by default, so the face has
% to be chosen explicitly.
%
% Latin Modern is the Computer Modern variety used here, and the choice is
% forced rather than aesthetic. Under T1 encoding, plain `cmr` has no Type 1
% outlines unless cm-super is installed, and it is not in the pinned
% distribution, so pdflatex falls back to EC bitmaps and every glyph in the
% PDF becomes a Type 3 font. That was measured: the first build produced a
% document with 10 Type 3 fonts and no Type 1. Latin Modern is Computer
% Modern redrawn with a full T1 repertoire and real .pfb outlines, so the
% same design comes out scalable.
\usepackage[T1]{fontenc}
\usepackage[utf8]{inputenc}
\usepackage{lmodern}
% microtype is not in the pinned TeX distribution. emergencystretch alone
% keeps the two-column measure from producing overfull lines.
\setlength{\emergencystretch}{3em}

\usepackage{amsmath}
\usepackage{amssymb}
\usepackage{amsthm}
\usepackage{booktabs}
\usepackage{array}
% enumitem is not in the pinned TeX distribution; the list spacing below
% does what \setlist would have done.
\usepackage{float}
\usepackage{xcolor}
\usepackage{listings}
\usepackage{tikz}
\usetikzlibrary{arrows.meta,positioning,fit,backgrounds,calc}
\usepackage[hidelinks,bookmarks=true]{hyperref}

% Numbered, definition-style box for a rule the reader is meant to apply.
\theoremstyle{definition}
\newtheorem{rulebox}{Rule}
\newtheorem{example}{Example}

% Unnumbered asides.
\theoremstyle{remark}
\newtheorem*{tip}{Tip}
\newtheorem*{pitfall}{Pitfall}

% Tighter lists, in place of enumitem's \setlist.
\setlength{\topsep}{2pt}
\setlength{\itemsep}{1pt}
\setlength{\parsep}{0pt}

% Code in the running text. The typewriter face under T1 ligates `<<`
% and `>>` into guillemets, so a shift printed as \code{<<} came out as
% a French quotation mark (issue 571). microtype's \DisableLigatures is
% not in the pinned distribution, but the primitive it rests on is
% pdfTeX's own: \pdfnoligatures strips every ligature from the font it
% is given, and the current font inside \texttt is the typewriter face
% at the size in use, so each size is stripped the first time \code
% sets it. Code wants no ligature at all: two quotes are two quotes.
\newcommand{\code}[1]{\texttt{\ifdefined\pdfnoligatures\pdfnoligatures\font\fi #1}}
\newcommand{\term}[1]{\emph{#1}}

% Syntax highlighting. Four roles, distinguishable in colour and still
% distinguishable in greyscale, because an IEEE article gets printed: the
% keyword is the only bold role, the comment is the only italic one, and the
% two remaining roles differ in lightness as well as in hue.
\definecolor{kwblue}{RGB}{20,60,150}
\definecolor{cmtgreen}{RGB}{90,120,90}
\definecolor{strred}{RGB}{150,50,40}
\definecolor{numgray}{gray}{0.45}
\definecolor{framegray}{gray}{0.70}
\definecolor{bgshade}{gray}{0.97}
% A darker fill for figure cells. The listing background has to stay very
% light to keep code readable; a figure cell has to be clearly shaded or the
% distinction it draws is invisible in print.
\definecolor{cellshade}{gray}{0.90}

% The listing size is the document's choice, because it depends on the
% column: a two-column article sets `\listingsize` to 6pt and keeps its
% listings within 70 columns, or to 5.6pt when it includes source that
% rustfmt sets at 80, which is what fits a column's frame (issue 1061);
% a one-column document keeps the body size and includes source files
% at 80. `//docs:frame_test` checks every listing line against its frame. Lines are never broken; a line that does not fit
% overflows the frame, which is visible, rather than wrapping, which is
% not.
\providecommand{\listingsize}{\scriptsize}

% `'` is deliberately not a string delimiter: the language uses it for loop
% labels, as Rust does, so treating it as a quote would swallow the rest of
% the line.
\lstdefinestyle{txhdl}{
  basicstyle=\ttfamily\listingsize,
  keywordstyle=\bfseries\color{kwblue},
  commentstyle=\itshape\color{cmtgreen},
  stringstyle=\color{strred},
  numberstyle=\tiny\color{numgray},
  morekeywords={mod,use,pub,const,type,struct,enum,trait,impl,fn,async,let,mut,
    match,if,else,loop,while,for,in,break,continue,return,as,await,
    module,bus,channel,transaction,protocol,pipeline,flow,seq,config,
    port,stage,pipe,instance,connect,bind,tag,domain,blackbox,foreign,
    property,role,signal,handler,true,false,
    sequence,interface,modport,implements,package,import,var,wire,func,proc,
    case,default,next,fork,branch,join,state,fsm},
  morecomment=[l]{//},
  morestring=[b]",
  showstringspaces=false,
  columns=fullflexible,
  keepspaces=true,
  breaklines=false,
  % Framed, so a listing is bounded rather than floating in the column.
  frame=single,
  rulecolor=\color{framegray},
  framesep=3pt,
  backgroundcolor=\color{bgshade},
  xleftmargin=3pt,
  xrightmargin=3pt,
  aboveskip=6pt,
  belowskip=6pt,
  % Every listing is numbered and captioned, so it can be referred to by
  % number. `caption` without `float` numbers the listing and leaves it
  % where it was written, which is what a listing in running prose needs.
  captionpos=b,
  abovecaptionskip=2pt,
  belowcaptionskip=6pt,
}

% Figures drawn as text. Framed the same way, and with the highlighting
% switched off, because a box-and-arrow diagram is not code and half its
% words turning blue reads as an accident. Setting a style adds keys rather
% than clearing the ones already set, so the three roles are neutralised by
% hand rather than by leaving them out. Program output is printed in this
% style too, and a netlist's assignment can run past the frame, so a line
% that does not fit is broken and the rest indented; source never is.
\lstdefinestyle{ascii}{
  basicstyle=\ttfamily\listingsize,
  keywordstyle=\ttfamily,
  commentstyle=\ttfamily,
  stringstyle=\ttfamily,
  columns=fullflexible,
  keepspaces=true,
  breaklines=true,
  breakindent=2em,
  frame=single,
  rulecolor=\color{framegray},
  framesep=3pt,
  backgroundcolor=\color{bgshade},
  xleftmargin=3pt,
  xrightmargin=3pt,
  aboveskip=6pt,
  belowskip=6pt,
}
\lstset{style=txhdl}

% House TikZ styles. `shorten` lives on the arrow styles rather than on each
% arrow, so every arrow in the document gets the gap, including ones added
% later. TikZ clips a path at a node's border, so without it an arrowhead
% butts against the box with no gap at all. Keep `node distance` well above
% twice the shorten, or the arrow shrinks to nothing.
\tikzset{
  box/.style={draw,rounded corners=2pt,align=center,inner sep=4pt,
              font=\footnotesize},
  boxf/.style={box,fill=bgshade},
  lbl/.style={font=\scriptsize\itshape,align=center},
  fl/.style={-{Stealth[length=4pt]},shorten >=2.5pt,shorten <=2.5pt},
  fld/.style={fl,dashed},
}


% The contents list: the class gives a section's number 2.75em, which
% a Roman numeral past thirty overruns onto the title, so the box is
% wider here, and a subsection's entry starts where a title does.
\makeatletter
\def\l@section#1#2{\addpenalty{\@secpenalty}\addvspace{1.0em plus 1pt}%
    \@tempdima 5.75em \begingroup \parindent \z@ \rightskip \@pnumwidth%
    \parfillskip-\@pnumwidth {\bfseries\leavevmode #1}\hfil\hbox to\@pnumwidth{\hss #2}\par%
    \endgroup}
\def\l@subsection{\@dottedtocline{2}{5.75em}{5em}}
\def\l@subsubsection{\@dottedtocline{3}{10.75em}{5.5em}}
\makeatother


\InputIfFileExists{buildstamp.tex}{}{}
\providecommand{\buildstamp}{Draft build (outside the Bazel flow).}

\title{TxHDL Documents}

\author{Filip Filmar and Dragiša Janković%
\thanks{Both authors are independent. Filip Filmar is the
corresponding author, at f@hdlfactory.com; Dragiša Janković is
reached at linkedin.com/in/dragisa-jankovic.}%
\thanks{This is the cover document of TxHDL. It names every document
the project produces, describes its role, and specifies where each is
built. When a document is added to the project, it is recorded here.}%
\thanks{The authors used a large language model, Claude, as an
assistant in exploring concepts and in drafting and constructing
the documents and implementations; every repository commit documents
this collaboration and includes the prompt sequence verbatim.}%
\thanks{\buildstamp}}

% Keep top-floated tables firmly at the page head.
\makeatletter
\setlength{\@fptop}{0pt}
\makeatother

\markboth{TxHDL Documents}{Filmar and Janković: TxHDL Documents}

\begin{document}
\maketitle

\section{The Language and Its Documents}

TxHDL is a hardware description language embedded directly in Rust. It began
as two independent experimental formalisms, LHDL and TxHDL, developed for the
same digital design project. Merged first under a unified Rust-shaped surface
syntax, the composite language was subsequently streamlined: every syntactic
construct already provided by the Rust host language was systematically
eliminated until only a clean domain library remained.

The documentation preserves this architectural trajectory. A reader seeking
the operational mechanics of the current language should begin with the trio
of core reports: the embedding rationale, the discrete-event runtime, and the
compiled reference examples. An engineer seeking to write synthesizable
hardware immediately should turn directly to the step-by-step tutorial.

Every document in this collection is built hermetically by Bazel from the
project repository, \code{HDL/txhdl}, against a pinned \LaTeX{} toolchain.
Executing \code{bazel build //docs/...} on any supported host reproduces the
entire corpus bit-for-bit. Release workflows validate every listing and diagram
nightly, publishing verified PDFs on release.

\section{The Documentation Map}

Tables~\ref{tab:lang_docs}, \ref{tab:hardware_docs}, and~\ref{tab:record_docs}
enumerate the complete corpus of project documents alongside the Bazel targets
that construct them.

Table~\ref{tab:lang_docs} presents the foundational reports detailing language
semantics, embedding theory, runtime mechanics, and introductory tutorials.
Table~\ref{tab:hardware_docs} surveys the complete digital systems implemented
in TxHDL, from modular AXI peripherals and network-on-chip fabrics to the
tapeout-ready RISC-V processor core. Table~\ref{tab:record_docs} collects the
overarching summaries, empirical codebase metrics, historical dynamics, and
the unified monograph.

\clearpage

\begin{table}[t]
\caption{Language foundations, runtime mechanics, and core tutorials.}
\label{tab:lang_docs}
\centering
\begin{tabular}{@{}p{46mm}p{92mm}l@{}}
\toprule
Document & Role and Scope & Target \\
\midrule
TxHDL: Unifying a Dataflow and a Transaction HDL
  & The architectural synthesis. Documents what each original language
    contributed, analyzes the four semantic conflicts between dataflow and
    transaction-level models, and explains their formal resolution. Serves as
    the design rationale for the language morphology.
  & \code{//docs:article} \\
\addlinespace
Deleting a Language: Embedding TxHDL in Rust
  & The definitive exposition of modern TxHDL. Details the rule that custom
    syntax is deleted in favour of standard Rust types, traits, and executors.
    Every claim regarding compiler behavior is verified through automated
    probes. Recommended starting point.
  & \code{//docs:embedding} \\
\addlinespace
The TxHDL Runtime
  & The operational engine. Incorporates the complete runtime source code
    verbatim from the tree, accompanied by an exact prose specification of its
    discrete-event execution model, cycle boundaries, and transactional state
    propagation.
  & \code{//docs:runtime} \\
\addlinespace
TxHDL by Example
  & A graded progression of compiled reference designs, from basic signal types
    to dual-clock asynchronous FIFOs, juxtaposing verified source listings with
    execution traces generated during the build.
  & \code{//docs:examples} \\
\addlinespace
TxHDL: A Hardware Description Language That Is a Rust Library
  & The comprehensive system paper. Details end-to-end architecture with
    complete schematics: host language embedding, discrete-event semantics,
    simulation traces, register-transfer lowering, and synthesis of the
    Vreteno core.
  & \code{//docs:paper} \\
\addlinespace
TxHDL, Step by Step
  & An eight-stage hands-on tutorial built around executable hardware modules:
    synchronous counters, the discrete-event model of time, control structures,
    inter-unit channels, named ports, multi-clock domains, waveform generation,
    and netlist validation.
  & \code{//docs:tutorial} \\
\addlinespace
TxHDL from Zero: A Blinky in Bazel
  & The bootstrap guide for an empty directory: installing Bazelisk and Git,
    configuring a minimal standalone workspace without checking out the full
    tree, and building, simulating, and lowering a hardware blinky across five
    concise configuration files.
  & \code{//docs:zero} \\
\addlinespace
Formal Verification in TxHDL
  & A practical guide to bounded model checking and safety property proofs.
    Proves the \code{ex\_formal} module across all input states using
    SymbiYosys, analyzes falsified assertions via counterexample traces, and
    reaches formal cover targets.
  & \code{//docs:prove} \\
\bottomrule
\end{tabular}
\end{table}

\section{The Working Notes}

Beneath \code{docs/} in the repository lie the foundational engineering notes
from which these publications were drafted. Preserved in raw Markdown, they are
intentionally retained alongside the formal papers: an architectural decision
survives only as long as its underlying trade-offs and alternatives remain
legible.

\begin{itemize}
\item \code{unification-analysis.md}: The initial comparative study of LHDL
  and original TxHDL, detailing feature overlaps, semantic divergences, and
  every defect uncovered in the legacy prototypes.
\item \code{syntax-decisions.md}: The rationale governing the surface syntax
  of the unified language, resolving each syntactic conflict in favour of
  idiomatic Rust conventions.
\item \code{rust-embedding.md}: The detailed experimental notes documenting
  the language reduction, recording compiler output from every type-system probe.
\item \code{document-build-plan.md}: The hermetic build specification for
  the documentation suite, explaining the rationale for vendoring \TeX{} packages
  and pinning the \LaTeX{} toolchain in Bazel.
\end{itemize}

\clearpage

\begin{table}[t]
\caption{Hardware subsystems, core architectures, and physical implementation.}
\label{tab:hardware_docs}
\centering
\begin{tabular}{@{}p{46mm}p{92mm}l@{}}
\toprule
Document & Role and Scope & Target \\
\midrule
A Reservation Station in TxHDL
  & Out-of-order execution primitives. Specifies the operational discipline of
    an instruction reservation station cycle by cycle, demonstrates
    parameterised multi-port dispatch with upstream queueing, and analyzes the
    lowered gate-level netlist.
  & \code{//docs:station} \\
\addlinespace
AXI in TxHDL
  & Standard on-chip interconnect protocols. Implements AXI4 point-to-point
    channels through two host transaction verbs and one peripheral verb,
    handling concurrent split-response transactions, outstanding request
    tracking, and structural routing.
  & \code{//docs:axi} \\
\addlinespace
Razboj: A Minimal GPU in TxHDL
  & The first complex peripheral on the bus. Implements a hardware rasteriser
    that consumes memory-mapped display lists and streams computed pixels
    across an AXI interface into a dedicated framebuffer, complete with
    rendered visual output and gate netlists.
  & \code{//docs:razboj} \\
\addlinespace
A Network on Chip in TxHDL
  & Scalable multi-agent interconnect. Implements a 2D mesh network-on-chip
    featuring dimension-order wormhole routing, dual virtual channels for
    deadlock avoidance, local router ports, and transactional packet-to-AXI
    translation bridges.
  & \code{//docs:noc} \\
\addlinespace
Ethernet in TxHDL
  & Media-access control on the system bus. Combines a dual-buffer full-frame
    MAC, a parallelised eight-step CRC-32 generator, and a three-register
    AXI-Lite control port, verified with loopback wire tests and on-board
    AX7A200B PHY operation.
  & \code{//docs:eth} \\
\addlinespace
HDMI in TxHDL
  & High-definition video generation. Implements video timing raster generators,
    an AXI-Lite addressable framebuffer, and an I2C control master that programs
    the physical transmitter, validated through whole-frame pixel tests and
    on-board SiI9134 output.
  & \code{//docs:hdmi} \\
\addlinespace
PCIe in TxHDL
  & High-speed host attachment. Couples an AXI master to a generated PCI Express
    endpoint core, exposing host-accessible memory and control registers mapped
    into BAR1, compiled through bitstream generation for the AX7A200B development
    board.
  & \code{//docs:pcie} \\
\addlinespace
The TxHDL Flagship
  & The complete system-on-chip demonstration bitstream. Integrates the Vreteno
    RISC-V core, system memory, Ethernet MAC, and HDMI rasterizer within a single
    design, coordinating three asynchronous clock domains, an SPI flash
    bootloader, and on-board peripherals.
  & \code{//docs:flagship} \\
\addlinespace
Vreteno: An RV32IMAC Core in TxHDL
  & The flagship processor implementation. An RV32IMAC RISC-V CPU verified cycle
    by cycle against an authoritative instruction-set reference model, accompanied
    by software test suites, execution waveforms, and physical Artix-7 FPGA
    deployment validating hardware equivalence.
  & \code{//docs:vreteno} \\
\addlinespace
A Tapeout Prep in TxHDL: Vreteno on an Open 45\,nm Node
  & Silicon implementation study. Takes the Vreteno core from RTL through an
    automated OpenROAD physical design flow targeting the FreePDK45
    standard-cell library, documenting floorplanning, cell placement,
    clock-tree synthesis, routing, timing closure, and tapeout readiness.
  & \code{//docs:tapeout} \\
\addlinespace
TxHDL Datasheets
  & Standardised component references. Standalone specification sheets for each
    reusable module in the library, detailing register maps, parameter
    interfaces, structural netlist hierarchies, temporal behavior, and
    automated regression test harnesses.
  & \code{//docs:datasheets} \\
\bottomrule
\end{tabular}
\end{table}

\clearpage

\begin{table}[t]
\caption{Overviews, codebase metrics, project dynamics, and the collected monograph.}
\label{tab:record_docs}
\centering
\begin{tabular}{@{}p{46mm}p{92mm}l@{}}
\toprule
Document & Role and Scope & Target \\
\midrule
TxHDL Documents
  & The master index and documentation roadmap. Outlines the corpus structure,
    details document roles and compilation targets, and establishes the
    triad verification rules governing the repository.
  & \code{//docs:cover} \\
\addlinespace
TxHDL: What Has Been Built
  & The executive summary. Synthesises the core architectural thesis, library
    ecosystem, RISC-V processor, raster GPU, and automated verification
    infrastructure into a concise multi-page survey of completed work and
    future directions.
  & \code{//docs:showcase} \\
\addlinespace
TxHDL Cheat Sheet
  & A dense single-page visual reference. Summarises core syntax, the
    discrete-event model of time, essential vocabulary, sample module anatomy,
    lowered netlists, and common build commands. Available as both PDF and PNG.
  & \begin{tabular}[t]{@{}l@{}}\code{//docs:cheatsheet}\\\code{//docs:cheatsheet\_png}\end{tabular} \\
\addlinespace
TxHDL by the Numbers
  & Quantitative repository archaeology. Measures source line counts across
    functional domains, programming languages, and individual files. Captures an
    immutable historical snapshot of project scale alongside reproducible
    measurement tooling.
  & \code{//docs:stats} \\
\addlinespace
The Dynamics of a Project
  & Chronological development history. Reconstructs engineering evolution from
    Git commit history, rendering dependency-constrained Gantt schedules across
    concurrent architectural initiatives and reflecting on lessons learned.
  & \code{//docs:dynamics} \\
\addlinespace
TxHDL: The Collected Monograph
  & The unified publication. Assembles every project paper and specification in
    canonical reading order into a single publication, \code{txhdl.pdf},
    complete with hierarchical PDF bookmarks for printing or comprehensive study.
  & \code{//docs:all} \\
\bottomrule
\end{tabular}
\end{table}

\section{How the Documents Are Kept True}

Three inviolable engineering rules ensure that this documentation remains in
perpetual synchronization with the underlying implementation. Enforced through
continuous integration and documented in \code{AGENTS.md}, every proposed change
is subject to their constraints:

\begin{enumerate}
\item \textbf{The Triad Rule:} A language feature exists only if it is
  implemented simultaneously in three locations: within the runtime library,
  in an automated compilation test or reference example, and in the relevant
  documentation. Omission in any single location renders the change incomplete.
\item \textbf{Verbatim Listing Inclusion:} Source code listings are never
  transcribed manually into \LaTeX{} sources. Every snippet is extracted
  directly from active tree files during the build. If a code example fails to
  compile or drifts from the implementation, the document build fails
  immediately.
\item \textbf{Empirical Probes as Evidence:} Claims regarding host language
  acceptance or rejection are backed by standalone compiler probes. Even
  negative outcomes are preserved as permanent tests, ensuring that compiler
  error messages serve as reproducible empirical evidence.
\end{enumerate}

\end{document}
